\documentclass[../m00-session-02.tex]{subfiles}
\begin{document}
\TopicStandaloneTitle{01}{Give data a readable shape}{Data foundations: from S3 objects to analysable tables}
\TopicDivider{01}{Give data a readable shape}{Distinguish storage from a trusted, documented data contract.}

\begin{frame}{What this topic establishes}
  \begin{LearningObjectives}{By the end of this topic, learners can}
    \begin{itemize}
      \item explain why an S3 prefix is not a database table; and
      \item state a simple rule for separating raw, curated and derived data.
    \end{itemize}
  \end{LearningObjectives}
\end{frame}

\begin{frame}{A data lake needs a readable convention}
  \centering
  \begin{tikzpicture}[node distance=1.0cm]
    \node[upgrad/data] (raw) {S3 / raw\\source-shaped};
    \node[upgrad/service,right=of raw] (curated) {S3 / curated\\typed + validated};
    \node[upgrad/model,right=of curated] (features) {S3 / features\\model-ready};
    \draw[upgrad/flow] (raw) -- node[above,font=\scriptsize,text=UpGradSlate]{clean} (curated);
    \draw[upgrad/flow] (curated) -- node[above,font=\scriptsize,text=UpGradSlate]{transform} (features);
  \end{tikzpicture}
  \vspace{.3cm}
  \begin{DefinitionBox}{Data contract}
    An agreed description of fields, types, owners, freshness and allowed values. It lets downstream work fail clearly instead of silently drifting.
  \end{DefinitionBox}
\end{frame}
\end{document}
